\par\medskip\noindent\textbf{The reduced support and its exact cohomology.}
Let $S$ denote the reduced support before the six balancing
operations. Label the Enriques component $M$, the component over
$R_0$ by $A\simeq\PP^1\times\PP^1$, the two Hopf components by
$Q_a,Q_b$, and the four exceptional ruled surfaces by
$E_1,\ldots,E_4\simeq\mathbb F_2$.
The actual quotient and resolution charts give
\[
\begin{array}{c|c|c}
\text{components}&\text{double curve}&\text{number}\\ \hline
M,A&D_0,D_\infty\simeq\PP^1&2\\
M,Q_j&J_j\text{ elliptic},\ j=a,b&2\\
A,E_i&Z_i\simeq\PP^1&4\\
M,E_i&U_{i0},U_{i\infty}\simeq\PP^1&8.
\end{array}                                                   \tag{C.EN.S1}
\]
There are precisely eight triple points
$D_\epsilon\cap Z_i\cap U_{i\epsilon}$, one for
$1\le i\le4$ and $\epsilon=0,\infty$, and no higher
intersection. On $A$, the curves $D_0,D_\infty$ are two
fibres of one projection and the $Z_i$ are sections of that
projection, all with the same divisor class. On $E_i$, $Z_i$
is the negative section and $U_{i0},U_{i\infty}$ are two
ruling fibres. These assertions retain the actual intersections:
the local quotient equation is $AC=B^2$, with base function
$T=u^2A$, and its $C$-chart after resolution has
$A=C\delta^2$, $B=C\delta$, hence $T=u^2C\delta^2$.
The three reduced components are $u=0$, $C=0$, and
$\delta=0$. Thus this support has simple normal crossings.

The elliptic map $M\to\PP^1$, induced by $z\mapsto z^2$,
has at $0,\infty$ the actual fibre divisors
\[
 \begin{gathered}
 F_0=2D_0+\sum_{i=1}^4U_{i0},\qquad
 F_\infty=2D_\infty+\sum_{i=1}^4U_{i\infty},\\
 [F_0]=[F_\infty]=2[J_a]=2[J_b]
       \quad\text{in }H_2(M,\Q).
 \end{gathered}                                          \tag{C.EN.S2}
\]
Indeed near a fixed base value the elliptic involution has
four fixed points; the displayed resolution gives one central
rational curve of multiplicity two and four exceptional arms
of multiplicity one. At the two root values the chart
$x=a^2b$ gives the two smooth double fibres $2J_a,2J_b$.
All these are fibres of the specified proper elliptic map,
which proves the equality of their rational fundamental classes.

The ten rational curves in (C.EN.S2) have precisely one rational
homology relation, namely $F_0-F_\infty=0$. To check this,
their intersection matrix is the direct sum of two affine
$D_4$ matrices: each central curve and its four arms has
self-intersection $-2$, each arm meets the central curve
once, and all remaining intersections are zero. Each matrix
has rank four and kernel generated by $(2,1,1,1,1)$.
Any homology relation lies in this two-dimensional kernel.
The common fibre class is nonzero, since its intersection
with an ample divisor of the projective surface $M$ is
positive. Thus the only kernel combination which is zero in
homology is the difference of the two fibres. In particular
these ten curves span a nine-dimensional subspace of
$H_2(M,\Q)$.

Here is the complete constant-sheaf comparison. If $S^{[p]}$
denotes the disjoint union of the $(p+1)$-fold intersections
listed in (C.EN.S1), the alternating restriction maps give
an exact complex
\[
 0\longrightarrow\Q_S\longrightarrow\Q_{S^{[0]}}
 \longrightarrow\Q_{S^{[1]}}
 \longrightarrow\Q_{S^{[2]}}\longrightarrow0.           \tag{C.EN.S3}
\]
At a point this is the augmented cochain complex of the
simplex formed by the branches through that point, proving
its exactness. Its cohomology spectral sequence has
$E_1^{p,q}=H^q(S^{[p]},\Q)$.
In degree zero the dimensions are $8,16,8$. The first
map has rank seven, because the incidence graph is connected.
The second has rank eight: the boundary of each triple
point contains its own edge $U_{i\epsilon}$, which occurs
in no other triangle. Consequently its degree-zero cohomology
has dimensions $(1,1,0)$ in columns $0,1,2$.

Only the two Hopf components contribute to $H^1(S^{[0]},\Q)$,
of total dimension two; only $J_a,J_b$ contribute to
$H^1(S^{[1]},\Q)$, of total dimension four.
Each Hopf restriction has rank one. In the actual presentation
$(\C^2\setminus0)/q^{2\Z}$, inclusion of its elliptic
attaching curve preserves the radial deck generator and
kills its angular generator. The two images lie in separate
elliptic summands. This map therefore has rank two.

The degree-two restriction map has dimensions $20\to16$
and rank fourteen. We verify its rank through the dual map
on fundamental curve classes. Give coefficients
$x_0,x_\infty$ to $D_0,D_\infty$, coefficients $y_a,y_b$
to $J_a,J_b$, coefficients $z_i$ to $Z_i$, and coefficients
$u_{i0},u_{i\infty}$ to $U_{i0},U_{i\infty}$, orienting
edges from the first component to the second in (C.EN.S1).
The $E_i$ equations force
$z_i=0$ and $u_{i\infty}=-u_{i0}$; the $A$ equation
forces $x_\infty=-x_0$. Equation (C.EN.S2) and the unique
rational-curve relation then force
\[
 y_a+y_b=0,\qquad u_{i0}=x_0/2\quad(1\le i\le4).
                                                               \tag{C.EN.S4}
\]
For completeness, the common half-fibre class is independent
of the four-dimensional span of the differences
$D_0-D_\infty,U_{i0}-U_{i\infty}$: in the curve vector
space modulo the sole relation $F_0-F_\infty$, a multiple
of $F_0$ cannot be a difference vector unless that multiple
is zero. Thus (C.EN.S4) lists all kernel elements. Its two
parameters are the relation $J_a-J_b$ and the relation
$2D_0-2D_\infty+\sum_i(U_{i0}-U_{i\infty})$.
This proves rank fourteen.

Finally the components contribute dimension two in degree
three, from the Hopf components, and dimension eight in
degree four. Intersections have no cohomology in those
degrees. The sole possible higher differential with nonzero
target would land in $E_2^{2,0}$, which was proved zero;
all groups $E_2^{2,q}$ for $q>0$ vanish because triple
intersections are points. The preceding actual maps give
\[
 \bigl(b_i(S)\bigr)_{i=0}^4=(1,1,8,4,8),\qquad e(S)=12.
                                                               \tag{C.EN.S5}
\]

The identical alternating complex with structure sheaves
is exact in the same normal crossing charts. All components
have $H^2(\OO)=0$; the only nonzero component $H^1(\OO)$
groups are those of the two Hopf surfaces, one-dimensional
each. Their restriction to $H^1(J_j,\OO_{J_j})$ is an
isomorphism. One can see this through the actual Hopf
elliptic fibration over $\PP^1$: its transition maps on the
fixed elliptic curve are translations, so its $R^1\OO$
is the trivial line; the Leray edge map is an isomorphism
and evaluation at any fibre is an isomorphism. The degree-zero
incidence calculation is unchanged. Hence
\[
 \bigl(h^i(S,\OO_S)\bigr)_{i=0}^2=(1,1,0).
                                                               \tag{C.EN.S6}
\]

The six negative balancing blocks may be placed away from
$M$, keeping the Enriques surface itself untouched. For
each point or curve blowup, the support is the full inverse
image. The relative-pair comparison and the restrictions
of the ambient exceptional Gysin classes give exactly
(C.KM.S2); each negative block adds $(2,6,2)$ in degrees
two, three, and four. For coherent cohomology the chosen
centre has generic reduced-support multiplicity one, and
the displayed blowup charts give $\sigma^*S=S'$ as
Cartier divisors, including the four crossing points of
the quartic step. Applying the projection formula to the
Cartier sequence, with $R\sigma_*\OO=\OO$, gives
$R\sigma_*\OO_{S'}=\OO_S$ as in (C.KM.S3).
Thus the balanced support $S_{\rm bal}$ has
\[
 \bigl(b_i(S_{\rm bal})\bigr)_{i=0}^4=(1,1,20,40,20),
 \qquad
 \bigl(h^i(S_{\rm bal},\OO_{S_{\rm bal}})\bigr)_{i=0}^2
       =(1,1,0),\qquad e(S_{\rm bal})=0.                 \tag{C.EN.S7}
\]
The complete incidence matrices, curve-class map, and
their exact ranks are executed in
\path{code/A_enriques_support.py} and retained in its
\texttt{.out} file.
